\documentclass[12pt,a4paper]{article}
\usepackage{fontspec}
\usepackage{xeCJK}
\setCJKmainfont{Microsoft JhengHei}
\setCJKsansfont{Microsoft JhengHei}
\setCJKmonofont{Microsoft JhengHei}
\usepackage{amsmath,amssymb,amsthm}
\usepackage{geometry}
\geometry{margin=2cm}
\usepackage{enumitem}
\usepackage{fancyhdr}
\usepackage{titlesec}
\usepackage{tcolorbox}
\tcbuselibrary{theorems,skins,breakable}
\usepackage{hyperref}
\hypersetup{colorlinks=true,linkcolor=blue!70!black,urlcolor=blue}

\pagestyle{fancy}
\fancyhf{}
\fancyhead[L]{\small 複變函數論期末複習}
\fancyhead[R]{\small \thepage}
\renewcommand{\headrulewidth}{0.4pt}

\newtcolorbox{theorembox}[1][]{
  colback=blue!5!white,
  colframe=blue!60!black,
  fonttitle=\bfseries,
  title=#1,
  breakable
}

\newtcolorbox{defbox}[1][]{
  colback=green!5!white,
  colframe=green!50!black,
  fonttitle=\bfseries,
  title=#1,
  breakable
}

\newtcolorbox{keybox}[1][]{
  colback=red!5!white,
  colframe=red!60!black,
  fonttitle=\bfseries,
  title=#1,
  breakable
}

\newtcolorbox{exbox}[1][]{
  colback=orange!5!white,
  colframe=orange!60!black,
  fonttitle=\bfseries,
  title=#1,
  breakable
}

\titleformat{\section}{\Large\bfseries}{第\thesection 章\quad}{0pt}{}
\titleformat{\subsection}{\large\bfseries}{\thesubsection\quad}{0pt}{}

\newcommand{\dd}{\,\mathrm{d}}
\newcommand{\ii}{\mathrm{i}}
\newcommand{\ee}{\mathrm{e}}
\newcommand{\Res}{\operatorname{Res}}
\newcommand{\PV}{\operatorname{P.V.}}

\begin{document}

\begin{center}
{\LARGE\bfseries 複變函數論\,---\,期末考複習筆記}\\[6pt]
{\large Chapter 5: 複數平面上的積分 \quad Chapter 6: 級數與留數}\\[4pt]
\today
\end{center}

\tableofcontents
\newpage

%=============================================================
\section{複數平面上的積分 (Integration in the Complex Plane)}
%=============================================================

%-------------------------------------------------------------
\subsection{5.1 實數積分回顧}
%-------------------------------------------------------------

\begin{defbox}[定積分]
$$\int_a^b f(x)\dd{x} = \lim_{\|P\|\to 0}\sum_{k=1}^{n}f(x_k^*)\,\Delta x_k = F(b)-F(a)$$
\end{defbox}

\subsubsection*{曲線的參數化}
曲線 $C$ 可藉由 $x=x(t),\ y=y(t),\ a\le t\le b$ 參數化，其中 $x(t),y(t)$ 為連續實函數。
\begin{itemize}[nosep]
\item \textbf{平滑曲線 (smooth curve)}：$z(t)$ 為一對一函數，$x'(t),y'(t)$ 連續且不同時為零。
\item \textbf{分段平滑曲線 (piecewise smooth curve)}：由有限個平滑曲線 $C_1,\dots,C_n$ 首尾相連。
\item \textbf{簡單曲線 (simple curve)}：曲線不自交（起終點除外）。
\item \textbf{封閉曲線 (closed curve)}：起點 $=$ 終點。
\item \textbf{簡單封閉曲線 (simple closed curve)}：不自交且封閉。
\end{itemize}

\subsubsection*{線積分的參數化計算}
\begin{align}
\int_C G(x,y)\dd{x} &= \int_a^b G\bigl(x(t),y(t)\bigr)\,\frac{\dd{x}}{\dd{t}}\dd{t} \\
\int_C G(x,y)\dd{y} &= \int_a^b G\bigl(x(t),y(t)\bigr)\,\frac{\dd{y}}{\dd{t}}\dd{t} \\
\int_C G(x,y)\dd{s} &= \int_a^b G\bigl(x(t),y(t)\bigr)\sqrt{\left(\frac{\dd{x}}{\dd{t}}\right)^2+\left(\frac{\dd{y}}{\dd{t}}\right)^2}\;\dd{t}
\end{align}

%-------------------------------------------------------------
\subsection{5.2 複數積分 (Complex Integrals)}
%-------------------------------------------------------------

\begin{defbox}[複數積分定義]
$$\int_C f(z)\dd{z} = \lim_{\|P\|\to 0}\sum_{k=1}^n f(z_k^*)\,\Delta z_k$$
\end{defbox}

\begin{itemize}[nosep]
\item $\oint_C f(z)\dd{z}$：逆時針為\textbf{正方向}（positively oriented）。
\item $\displaystyle\oint_{-C} f(z)\dd{z}$：順時針為\textbf{負方向}。
\end{itemize}

\begin{theorembox}[Theorem 5.2.1 --- 圍道積分的計算]
若 $f$ 在光滑曲線 $C: z(t)=x(t)+\ii y(t),\ a\le t\le b$ 上連續，則
$$\int_C f(z)\dd{z} = \int_a^b f\bigl(z(t)\bigr)\,z'(t)\dd{t}$$
\end{theorembox}

\begin{keybox}[計算步驟]
\begin{enumerate}[nosep]
\item 將 $z$ 參數化：$z=z(t)$，求 $\dfrac{\dd{z}}{\dd{t}}$。
\item 將 $f(z)$ 以 $t$ 表示。
\item 計算 $\displaystyle\int_a^b f(z(t))\cdot z'(t)\dd{t}$。
\end{enumerate}
\end{keybox}

\begin{exbox}[典型範例]
計算 $\displaystyle\oint_C \frac{1}{z}\dd{z}$，其中 $C:|z|=1$ 逆時針。

令 $z=\ee^{\ii t},\ 0\le t\le 2\pi$，則 $\dd{z}=\ii\ee^{\ii t}\dd{t}$：
$$\oint_C \frac{1}{z}\dd{z}=\int_0^{2\pi}\frac{1}{\ee^{\ii t}}\cdot\ii\ee^{\ii t}\dd{t}=\int_0^{2\pi}\ii\dd{t}=2\pi\ii$$
\end{exbox}

\begin{theorembox}[Theorem 5.2.2 --- 圍道積分性質]
\begin{enumerate}[nosep]
\item $\displaystyle\int_C kf(z)\dd{z}=k\int_C f(z)\dd{z}$
\item $\displaystyle\int_C[f(z)+g(z)]\dd{z}=\int_C f(z)\dd{z}+\int_C g(z)\dd{z}$
\item $\displaystyle\int_C f(z)\dd{z}=\int_{C_1}f(z)\dd{z}+\int_{C_2}f(z)\dd{z}$（$C=C_1+C_2$ 首尾相連）
\item $\displaystyle\int_{-C}f(z)\dd{z}=-\int_C f(z)\dd{z}$（反向）
\end{enumerate}
\end{theorembox}

\begin{theorembox}[Theorem 5.2.3 --- ML 不等式（估界定理）]
若 $f$ 在光滑曲線 $C$ 上連續且 $|f(z)|\le M$，則
$$\left|\int_C f(z)\dd{z}\right|\le ML$$
其中 $L$ 為曲線 $C$ 的長度，$L=\displaystyle\int_a^b\sqrt{[x'(t)]^2+[y'(t)]^2}\dd{t}$。
\end{theorembox}

\begin{exbox}[ML 不等式範例]
求 $\displaystyle\left|\oint_C\frac{\ee^z}{z+1}\dd{z}\right|$ 的上界，$C:|z|=4$。

$L=2\pi\cdot4=8\pi$。在 $|z|=4$ 上：$|\ee^z|\le \ee^4$，$|z+1|\ge||z|-1|=3$。

故 $M=\dfrac{\ee^4}{3}$，上界 $=ML=\dfrac{8\pi\ee^4}{3}$。
\end{exbox}

%-------------------------------------------------------------
\subsection{5.3 Cauchy--Goursat 定理}
%-------------------------------------------------------------

\begin{defbox}[連通區域]
\begin{itemize}[nosep]
\item \textbf{單連通區域 (simply connected domain)}：任何簡單封閉曲線都可連續收縮為一點，仍保持在該 domain 內。（無洞）
\item \textbf{多連通區域 (multiply connected domain)}：有一個洞為 doubly connected；兩個洞為 triply connected。
\end{itemize}
\end{defbox}

\begin{theorembox}[Theorem 5.3.1 --- Cauchy--Goursat 定理]
若 $f$ 在\textbf{單連通區域} $D$ 內解析，則對 $D$ 內任意簡單封閉曲線 $C$：
$$\oint_C f(z)\dd{z}=0$$
\end{theorembox}

\begin{keybox}[判斷要點]
\begin{itemize}[nosep]
\item 確認 $f$ 在 $C$ 所圍區域內\textbf{處處解析}（無奇異點）。
\item 若不可解析的點在 $C$ 以外，積分仍為 $0$。
\item 若不可解析的點在 $C$ 以內，\textbf{不能}直接用此定理。
\end{itemize}
\end{keybox}

\begin{theorembox}[Theorem 5.3.2 --- 多連通區域的 Cauchy--Goursat 定理]
若 $C,C_1,\dots,C_n$ 為正向簡單封閉曲線，$C_1,\dots,C_n$ 在 $C$ 內部且互不交，$f$ 在 $C$ 與所有 $C_k$ 之間的區域解析，則：
$$\oint_C f(z)\dd{z}=\sum_{k=1}^n\oint_{C_k}f(z)\dd{z}$$
\end{theorembox}

\begin{keybox}[路徑變形原理]
若 $C$ 包含非解析點（doubly connected domain），可重新定義 $C$ 為包含該點的更小圓 $C_1$，利用
$$\oint_C f(z)\dd{z}=\oint_{C_1}f(z)\dd{z}$$
$C_1$ 可選為方便計算的圓，半徑任意。
\end{keybox}

\begin{exbox}[重要結果]
$$\oint_C\frac{\dd{z}}{(z-z_0)^n}=\begin{cases}2\pi\ii, & n=1\\ 0, & n\ge 2 \text{ 或 } n\le 0\end{cases}$$
其中 $C$ 為包含 $z_0$ 的任意簡單封閉曲線。
\end{exbox}

%-------------------------------------------------------------
\subsection{5.4 路徑獨立性 (Independence of Path)}
%-------------------------------------------------------------

\begin{theorembox}[Theorem 5.4.1 --- 解析性蘊含路徑獨立]
若 $f$ 在\textbf{單連通區域} $D$ 內解析，$C$ 為 $D$ 內任意 contour，則 $\displaystyle\int_C f(z)\dd{z}$ 與路徑無關。
\end{theorembox}

\begin{keybox}[應用]
若 $f$ 為 entire function（整函數，處處解析），則可自由選擇路徑計算積分。例如 $2z$ 為 entire function，可選擇任何方便的路徑。
\end{keybox}

%-------------------------------------------------------------
\subsection{5.5 Cauchy 積分公式（若考試範圍包含）}
%-------------------------------------------------------------

\begin{theorembox}[Cauchy 積分公式]
若 $f$ 在單連通區域 $D$ 內解析，$z_0$ 為 $D$ 中一點，$C$ 為 $D$ 內包含 $z_0$ 的正向簡單封閉曲線，則：
$$f(z_0)=\frac{1}{2\pi\ii}\oint_C\frac{f(z)}{z-z_0}\dd{z}$$
推廣：
$$f^{(n)}(z_0)=\frac{n!}{2\pi\ii}\oint_C\frac{f(z)}{(z-z_0)^{n+1}}\dd{z}$$
\end{theorembox}

\newpage
%=============================================================
\section{級數與留數 (Series and Residues)}
%=============================================================

%-------------------------------------------------------------
\subsection{6.1 序列與級數}
%-------------------------------------------------------------

\begin{defbox}[序列收斂]
序列 $\{z_n\}$ 收斂至 $L=a+\ii b$ $\iff$ $\text{Re}(z_n)\to a$ 且 $\text{Im}(z_n)\to b$。
\end{defbox}

\subsubsection*{級數 (Series)}
$$\sum_{k=1}^{\infty}z_k = z_1+z_2+z_3+\cdots$$

\begin{theorembox}[幾何級數]
$$\sum_{k=1}^{\infty}az^{k-1}=\frac{a}{1-z},\quad |z|<1$$
\end{theorembox}

\begin{keybox}[特殊幾何級數（必背！）]
\begin{align}
\frac{1}{1-z}&=1+z+z^2+z^3+\cdots, \quad |z|<1 \\[4pt]
\frac{1}{1+z}&=1-z+z^2-z^3+\cdots, \quad |z|<1
\end{align}
\end{keybox}

\begin{theorembox}[Theorem 6.1.2 --- 收斂的必要條件]
若 $\displaystyle\sum_{k=1}^{\infty}z_k$ 收斂，則 $\displaystyle\lim_{n\to\infty}z_n=0$。
\end{theorembox}

\begin{theorembox}[Theorem 6.1.3 --- 第 $n$ 項發散判別法]
若 $\displaystyle\lim_{n\to\infty}z_n\ne 0$，則 $\displaystyle\sum_{k=1}^{\infty}z_k$ 發散。
\end{theorembox}

\begin{defbox}[絕對收斂與條件收斂]
\begin{itemize}[nosep]
\item \textbf{絕對收斂}：$\displaystyle\sum_{k=1}^{\infty}|z_k|$ 收斂 $\Rightarrow$ $\displaystyle\sum_{k=1}^{\infty}z_k$ 絕對收斂。
\item \textbf{條件收斂}：$\displaystyle\sum z_k$ 收斂但 $\displaystyle\sum|z_k|$ 發散。
\end{itemize}
\end{defbox}

\subsubsection*{$p$-級數}
$$\sum_{k=1}^{\infty}\frac{1}{k^p}\quad\text{收斂若 }p>1;\quad \text{發散若 }p\le 1.$$

\begin{theorembox}[Theorem 6.1.4 --- 比值審斂法 (Ratio Test)]
設 $\displaystyle L=\lim_{n\to\infty}\left|\frac{z_{n+1}}{z_n}\right|$，則：
\begin{itemize}[nosep]
\item $L<1$：絕對收斂
\item $L>1$ 或 $L=\infty$：發散
\item $L=1$：無法判定
\end{itemize}
\end{theorembox}

\begin{theorembox}[Theorem 6.1.5 --- 根式審斂法 (Root Test)]
設 $\displaystyle L=\lim_{n\to\infty}\sqrt[n]{|z_n|}$，則：
\begin{itemize}[nosep]
\item $L<1$：絕對收斂
\item $L>1$ 或 $L=\infty$：發散
\item $L=1$：無法判定
\end{itemize}
\end{theorembox}

%-------------------------------------------------------------
\subsection{6.1 續：冪級數 (Power Series)}
%-------------------------------------------------------------

\begin{defbox}[冪級數]
以 $z_0$ 為中心：
$$\sum_{k=0}^{\infty}a_k(z-z_0)^k = a_0+a_1(z-z_0)+a_2(z-z_0)^2+\cdots$$
其中 $a_k$ 為複數常數。
\end{defbox}

\begin{keybox}[收斂半徑 $R$]
冪級數 $\displaystyle\sum_{k=0}^{\infty}a_k(z-z_0)^k$ 在 $|z-z_0|<R$ 收斂，$|z-z_0|>R$ 發散。
\begin{itemize}[nosep]
\item $R=0$：僅在 $z=z_0$ 收斂（其餘發散）
\item $0<R<\infty$：$|z-z_0|=R$ 為收斂圓
\item $R=\infty$：在整個複數平面上收斂
\end{itemize}

\textbf{求 $R$ 的方法}：用 Ratio Test 或 Root Test，令結果 $<1$，解出 $|z-z_0|<R$。
\end{keybox}

\begin{exbox}[收斂半徑範例]
$\displaystyle\sum_{k=1}^{\infty}\frac{(-1)^{k+1}}{k!}(z-1-\ii)^k$，中心 $z_0=1+\ii$。

Ratio Test：$\displaystyle\lim_{n\to\infty}\left|\frac{z-1-\ii}{n+1}\right|=0<1$ 對所有 $z$ 成立。

故 $R=\infty$。
\end{exbox}

%-------------------------------------------------------------
\subsection{6.2 Taylor 級數}
%-------------------------------------------------------------

\begin{keybox}[冪級數在收斂半徑內的性質]
在 $|z-z_0|<R$ 內，冪級數 $\displaystyle\sum_{k=0}^{\infty}a_k(z-z_0)^k$：
\begin{enumerate}[nosep]
\item \textbf{連續} (Thm 6.2.1)
\item \textbf{可逐項微分} (Thm 6.2.2)
\item \textbf{可逐項積分} (Thm 6.2.3)
\end{enumerate}
\end{keybox}

\begin{theorembox}[Theorem 6.2.4 --- Taylor 定理]
若 $f$ 在區域 $D$ 內解析，$z_0\in D$，則 $f$ 有級數表示：
$$f(z)=\sum_{k=0}^{\infty}\frac{f^{(k)}(z_0)}{k!}(z-z_0)^k$$
對以 $z_0$ 為中心、完全在 $D$ 內的最大圓有效。
\end{theorembox}

\begin{keybox}[Taylor 級數的係數]
$$a_k=\frac{f^{(k)}(z_0)}{k!}, \quad k=0,1,2,\dots$$
即 $a_0=f(z_0),\; a_1=f'(z_0),\; a_2=\dfrac{f''(z_0)}{2!},\;\dots$
\end{keybox}

\begin{keybox}[收斂半徑 $R$ = 中心 $z_0$ 到最近 isolated singularity 的距離]
\end{keybox}

\subsubsection*{Maclaurin 級數（$z_0=0$ 的 Taylor 級數）---\,必背！}

\begin{keybox}[三大 Maclaurin 級數]
\begin{align}
\ee^z &= 1+z+\frac{z^2}{2!}+\frac{z^3}{3!}+\cdots = \sum_{k=0}^{\infty}\frac{z^k}{k!}, \quad R=\infty \\[6pt]
\sin z &= z-\frac{z^3}{3!}+\frac{z^5}{5!}-\cdots = \sum_{k=0}^{\infty}(-1)^k\frac{z^{2k+1}}{(2k+1)!}, \quad R=\infty \\[6pt]
\cos z &= 1-\frac{z^2}{2!}+\frac{z^4}{4!}-\cdots = \sum_{k=0}^{\infty}(-1)^k\frac{z^{2k}}{(2k)!}, \quad R=\infty
\end{align}
\end{keybox}

\begin{exbox}[Taylor 展開技巧]
\textbf{方法一}：直接利用已知 Maclaurin 級數做代換。\\
例：$\ee^{-z^2}$：將 $z$ 換成 $-z^2$ 代入 $\ee^z$ 的展開式。\\
例：$\sin 3z$：將 $z$ 換成 $3z$ 代入 $\sin z$ 的展開式。

\textbf{方法二}：利用微分或積分。\\
例：$\dfrac{1}{(1-z)^2}=\left(\dfrac{1}{1-z}\right)'=\displaystyle\sum_{k=1}^{\infty}kz^{k-1}$，$|z|<1$。

\textbf{方法三}：改寫成 $\dfrac{a}{1\pm\square}$ 的形式再展開。\\
例：$f(z)=\dfrac{1}{1-z}$ 以 $z_0=2\ii$ 為中心：
$$f(z)=\frac{1}{1-2\ii-(z-2\ii)}=\frac{1}{1-2\ii}\cdot\frac{1}{1-\frac{z-2\ii}{1-2\ii}}$$
再用 $\dfrac{1}{1-w}=\displaystyle\sum_{k=0}^{\infty}w^k$ 展開。
\end{exbox}

%-------------------------------------------------------------
\subsection{6.3 Laurent 級數}
%-------------------------------------------------------------

\begin{defbox}[奇異點 (Singularity)]
\begin{itemize}[nosep]
\item 若 $f(z)$ 在 $z=z_0$ 不解析，則 $z_0$ 為 $f$ 的\textbf{奇異點}（singular point）。
\item \textbf{孤立奇異點 (isolated singularity)}：存在去心鄰域使 $f$ 在其中解析。
\item \textbf{非孤立奇異點}：例如 $f(z)=\ln z$，分支切線上每一點。
\end{itemize}
\end{defbox}

\begin{theorembox}[Theorem 6.3.1 --- Laurent 定理]
若 $f$ 在環形區域 $r<|z-z_0|<R$ 內解析，則
$$f(z)=\sum_{k=-\infty}^{\infty}a_k(z-z_0)^k$$
其中
$$a_k=\frac{1}{2\pi\ii}\oint_C\frac{f(z)}{(z-z_0)^{k+1}}\dd{z},\quad k=0,\pm1,\pm2,\dots$$
$C$ 為環形區域內的任意正向簡單封閉曲線。
\end{theorembox}

\begin{keybox}[Laurent 級數的結構]
$$f(z)=\underbrace{\cdots+\frac{a_{-2}}{(z-z_0)^2}+\frac{a_{-1}}{z-z_0}}_{\text{主部 (principal part)}}+\underbrace{a_0+a_1(z-z_0)+a_2(z-z_0)^2+\cdots}_{\text{解析部分 (analytic part)}}$$
\end{keybox}

\begin{keybox}[Laurent 展開的核心技巧]
\textbf{目標}：將 $f(z)$ 寫成 $\dfrac{a}{1\pm\square}$ 的形式，其中 $|\square|<1$。

根據環形區域的條件判斷哪個量小於 1：
\begin{itemize}[nosep]
\item 若 $|z|<R$（$z$ 在小圓內）：令 $\square$ 中含 $z$（或 $z-z_0$），使 $|\square|<1$。
\item 若 $|z|>r$（$z$ 在大圓外）：令 $\square$ 中含 $\dfrac{1}{z}$（或 $\dfrac{1}{z-z_0}$），使 $|\square|<1$。
\end{itemize}
\end{keybox}

\begin{exbox}[經典範例：$f(z)=\dfrac{1}{z(z-1)}$ 的四種 Laurent 展開]
\textbf{(a) $0<|z|<1$：}
$$f(z)=\frac{-1}{z}\cdot\frac{1}{1-z}=\frac{-1}{z}(1+z+z^2+\cdots)=\frac{-1}{z}-1-z-z^2-\cdots$$

\textbf{(b) $1<|z|$（即 $|z|>1$）：}
$$f(z)=\frac{1}{z^2}\cdot\frac{1}{1-\frac{1}{z}}=\frac{1}{z^2}\left(1+\frac{1}{z}+\frac{1}{z^2}+\cdots\right)=\frac{1}{z^2}+\frac{1}{z^3}+\frac{1}{z^4}+\cdots$$
（全部為主部）

\textbf{(c) $0<|z-1|<1$：}
$$f(z)=\frac{1}{z-1}\cdot\frac{1}{1+(z-1)}=\frac{1}{z-1}\bigl[1-(z-1)+(z-1)^2-\cdots\bigr]$$

\textbf{(d) $1<|z-1|$：}
$$f(z)=\frac{1}{(z-1)^2}\cdot\frac{1}{1+\frac{1}{z-1}}=\frac{1}{(z-1)^2}-\frac{1}{(z-1)^3}+\cdots$$
\end{exbox}

\begin{exbox}[利用已知級數做 Laurent 展開]
$f(z)=\ee^{3/z}$，$0<|z|<\infty$。

將 $\ee^w=\displaystyle\sum_{k=0}^{\infty}\frac{w^k}{k!}$ 中 $w=\dfrac{3}{z}$ 代入：
$$\ee^{3/z}=\sum_{k=0}^{\infty}\frac{(3/z)^k}{k!}=1+\frac{3}{z}+\frac{9}{2!\,z^2}+\frac{27}{3!\,z^3}+\cdots$$
由 Ratio Test：只要 $z\ne 0$ 即收斂。
\end{exbox}

\begin{exbox}[展開 $\dfrac{\sin z}{z^4}$ 做 Laurent 級數]
$$\frac{\sin z}{z^4}=\frac{1}{z^4}\left(z-\frac{z^3}{3!}+\frac{z^5}{5!}-\cdots\right)=\frac{1}{z^3}-\frac{1}{3!\,z}+\frac{z}{5!}-\frac{z^3}{7!}+\cdots$$
主部為 $\dfrac{1}{z^3}-\dfrac{1}{6z}$，解析部分為 $\dfrac{z}{5!}-\cdots$。有效範圍 $|z|>0$。
\end{exbox}

\newpage
%=============================================================
\section{考試常用公式與技巧速查}
%=============================================================

\subsection*{A. 幾何級數公式}
\begin{align}
\frac{1}{1-w}&=\sum_{k=0}^{\infty}w^k, \quad |w|<1 \\[4pt]
\frac{1}{1+w}&=\sum_{k=0}^{\infty}(-1)^k w^k, \quad |w|<1
\end{align}

\subsection*{B. 三大 Maclaurin 展開}
\begin{align}
\ee^z&=\sum_{k=0}^{\infty}\frac{z^k}{k!} & \sin z&=\sum_{k=0}^{\infty}(-1)^k\frac{z^{2k+1}}{(2k+1)!} & \cos z&=\sum_{k=0}^{\infty}(-1)^k\frac{z^{2k}}{(2k)!}
\end{align}

\subsection*{C. 圍道積分關鍵結果}
$$\oint_C\frac{\dd{z}}{(z-z_0)^n}=\begin{cases}2\pi\ii & n=1\\0 & n\ne 1\end{cases}$$
（$C$ 包含 $z_0$，正向）

\subsection*{D. Cauchy--Goursat 定理速查}
\begin{center}
\begin{tabular}{|l|c|}
\hline
\textbf{條件} & \textbf{結果} \\
\hline
$f$ 在 $C$ 內部解析（simply connected） & $\oint_C f\dd{z}=0$ \\
$C$ 不包含奇異點 & $\oint_C f\dd{z}=0$ \\
$C$ 包含奇異點（multiply connected） & 用路徑變形或部分分式 \\
\hline
\end{tabular}
\end{center}

\subsection*{E. 收斂性判別流程}

\begin{enumerate}[nosep]
\item 先用\textbf{第 $n$ 項判別法}：若 $\lim z_n\ne 0$，直接發散。
\item 若含階乘 $k!$ 或指數，用 \textbf{Ratio Test}。
\item 若第 $k$ 項整體是 $(\cdots)^k$ 形式，用 \textbf{Root Test}。
\item 可比較 $p$-級數：$\sum 1/k^p$，$p>1$ 收斂。
\end{enumerate}

\subsection*{F. Laurent 展開策略}

\begin{enumerate}[nosep]
\item 做\textbf{部分分式}分解。
\item 根據環形區域決定哪個量 $<1$。
\item 改寫成 $\dfrac{a}{1\pm\square}$ 形式（$|\square|<1$）。
\item 展開為幾何級數。
\item 或利用已知 Maclaurin 級數做代換（$\ee^{1/z},\;\dfrac{\sin z}{z^n}$ 等）。
\end{enumerate}

\subsection*{G. Taylor 級數收斂半徑的求法}

\begin{enumerate}[nosep]
\item \textbf{Ratio Test}：$\displaystyle\lim_{n\to\infty}\left|\frac{a_{n+1}}{a_n}(z-z_0)\right|<1$，解出 $|z-z_0|<R$。
\item \textbf{Root Test}：$\displaystyle\lim_{n\to\infty}\sqrt[n]{|a_n|}\cdot|z-z_0|<1$。
\item \textbf{幾何意義}：$R=$ 中心 $z_0$ 到 $f$ 最近\textbf{孤立奇異點}的距離。
\end{enumerate}

\end{document}
